\section{Torsion in a fixed complex dimension}
\label{geo:section}
\label{sec:geometry}

For every integer $m\geq1$, we will construct a simply connected closed
K\"ahler manifold $X_m$ of complex dimension nine by blowing up $\CP^9$
along $2m$ disjoint centers. Two quotient surfaces supply torsion at
primes two and three. A simultaneous projection places arbitrarily many
copies of the centers in $\CP^9$, and the integral blow-up formula then
records their contributions to homology. The surface construction and
blow-up argument also appear in \cite[Section~3]{OpenAICompanion}; the
projection step is what keeps the dimension independent of $m$.

\subsection{Two projective quotient surfaces}

Let
\[
 E=\C/(\Z+i\Z),\qquad
 \zeta_p=\exp(2\pi i/p),\qquad
 \Lambda_2=\Z+i\Z,\qquad \Lambda_3=\Z+\zeta_3\Z,
\]
and set $E_p=\C/\Lambda_p$ for $p\in\{2,3\}$. The automorphism
\[
 \sigma_p:E\times E_p\longrightarrow E\times E_p,
 \qquad \sigma_p(z,w)=(z+1/p,\zeta_p w)
\]
is well-defined because multiplication by $\zeta_p$ preserves
$\Lambda_p$. It has order $p$. Every nonidentity power translates the
first factor by $k/p$ for some $0<k<p$, and hence has no fixed point.
Thus
\[
 S_p=(E\times E_p)/\langle\sigma_p\rangle
\]
is a smooth compact connected complex surface. These are examples of
bielliptic surfaces, whose integral topology is studied in
\cite[Theorem~4.3]{Serrano1991}. The quotient presentations and homology
groups used here are recorded in \cite[Sections~2.2--2.3]{Blomme2025}.
We give the projectivity and homology calculations needed for the
construction.

\begin{lemma}
\label{geo:surfaces}
The surfaces $S_2$ and $S_3$ are projective. With
\[
 T_2=(\Z/2)^2,\qquad T_3=\Z/3,
\]
their integral homology groups are
\begin{equation}
\label{geo:surface-homology}
 H_i(S_p;\Z)\cong
 \begin{cases}
  \Z,&i=0,4,\\
  \Z^2\oplus T_p,&i=1,2,\\
  \Z^2,&i=3,\\
  0,&\text{otherwise}.
 \end{cases}
\end{equation}
In particular, each surface has total Betti number eight.
\end{lemma}

\begin{proof}
We first realize the quotient as a projective submanifold. A line bundle
of degree three on an elliptic curve has three independent sections.
Riemann--Roch applied after subtracting one or two points, allowing
multiplicity, shows that its complete linear system has no base points
and separates points and tangent directions. It therefore embeds the
curve in $\CP^2$. Taking the two curve embeddings and their Segre product
gives an embedding
\[
 \iota_p:E\times E_p\hookrightarrow\CP^8.
\]

A point of $\CP^8$ represents a linear polynomial on $(\C^9)^*$, up to
a nonzero scalar. Associate to an orbit the product of the $p$ linear
polynomials representing
\[
 \iota_p(x),\ \iota_p(\sigma_p x),\ \ldots,
 \ \iota_p(\sigma_p^{p-1}x).
\]
This gives a well-defined holomorphic map
\[
 S_p\longrightarrow\mathbb P(\operatorname{Sym}^p\C^9).
\]
Unique factorization shows that equal projective products have the same
unordered factors. Thus the map separates the orbits.

To check immersion, choose local representatives $L_0,\ldots,L_{p-1}$
of the orbit points. A tangent vector in the kernel of the derivative
would give
\[
 \sum_{j=0}^{p-1}\dot L_j\prod_{k\ne j}L_k
       =a\prod_{k=0}^{p-1}L_k
\]
for a scalar $a$. The factors are pairwise nonproportional because the
action is free. Restricting this identity to the hyperplane $L_j=0$
forces $\dot L_j$ to vanish there: the product of the remaining factors
is a nonzero polynomial on that hyperplane. Consequently $\dot L_j$ is
a scalar multiple of $L_j$ for every $j$. In particular, the projective
derivative of $\iota_p(x)$ is zero, so the original tangent vector is
zero. The quotient map is a local biholomorphism, and therefore the
induced map on $S_p$ is an injective holomorphic immersion. Compactness
makes it an embedding. This proves projectivity.

For the homology, the universal cover $\C^2\to S_p$ has deck group
generated by
\[
 a(z,w)=(z+1/p,\zeta_p w),\quad
 d(z,w)=(z+i,w),\quad
 t_\lambda(z,w)=(z,w+\lambda)\quad(\lambda\in\Lambda_p).
\]
Here $a^p$ is translation by $1$ in the first coordinate. The elements
$a,d$ commute and generate a free abelian group of rank two. Their
conjugation actions on $\Lambda_p$ are multiplication by $\zeta_p$ and
the identity, respectively. The deck group is therefore
$\Lambda_p\rtimes\Z^2$, and its abelianization is
\begin{equation}
\label{geo:abelianization}
 H_1(S_p;\Z)\cong
 \Z^2\oplus\Lambda_p/(\zeta_p-1)\Lambda_p.
\end{equation}
For $p=2$, multiplication by $\zeta_p-1$ is $-2I$. For $p=3$, its
matrix in the basis $(1,\zeta_3)$ is
\[
 \begin{pmatrix}-1&-1\\1&-2\end{pmatrix}.
\]
This matrix has determinant three and relatively prime entries, so its
Smith normal form has diagonal $(1,3)$. The finite summands in
\eqref{geo:abelianization} are exactly $T_2$ and $T_3$.

The complex orientation makes $S_p$ a closed oriented real
four-manifold. Integral Poincar\'e duality and the universal coefficient
theorem imply
\[
 \Tor H_2(S_p;\Z)\cong\Tor H^2(S_p;\Z)
       \cong\Tor H_1(S_p;\Z)=T_p,
 \qquad H_3(S_p;\Z)\cong\Z^2.
\]
Finally, the degree-$p$ covering gives
$p\chi(S_p)=\chi(E\times E_p)=0$. The ranks already computed in degrees
$0,1,3,4$ are $1,2,2,1$, so the rank in degree two is two. This proves
\eqref{geo:surface-homology}.
\end{proof}

Define the two types of blow-up center by
\begin{equation}
\label{geo:centers}
 Z_2=S_2,\qquad Z_3=S_3\times\CP^1\times\CP^1.
\end{equation}
They are smooth connected projective varieties of complex dimensions
two and four; projectivity of $Z_3$ follows from the Segre embedding.
The projective-line factors distribute the $3$-primary torsion among
shifts by $0,2,4$, with multiplicities $1,2,1$. Together with the
different blow-up codimensions, this will make which prime has the larger
multiplicity depend on the degree; Section~\ref{sec:counts}
computes the resulting generator counts.

\subsection{Simultaneous projection of the centers}

The next lemma shows that the number of centers does not force the
ambient dimension to increase. We use the classical secant-and-tangent
projection argument; compare \cite[Chapter~IV, Exercise~3.11(a)]{Hartshorne1977}.
For a disjoint union, secants joining different components must be
avoided as well as secants within one component.

\begin{lemma}[Simultaneous projective embedding]
\label{geo:projection}
Let $d\geq0$ be an integer, and let $V_1,\ldots,V_s$ be a finite family
of smooth complex projective varieties of complex dimension at most
$d$. Their disjoint union admits a holomorphic embedding into
$\CP^{2d+1}$.
\end{lemma}

\begin{proof}
Splitting disconnected varieties into their finitely many connected
components, we may assume each $V_\alpha$ is connected.
Choose projective embeddings $V_\alpha\hookrightarrow\mathbb P(W_\alpha)$
and place their images in the pairwise disjoint linear subspaces
$\mathbb P(W_\alpha)\subset\mathbb P(\bigoplus_\alpha W_\alpha)$.
Adding unused coordinates if necessary, we obtain a disjoint embedded
union $U\subset\CP^D$ with $D>2d+1$. We prove that projection from a
suitable point of $\CP^D$ embeds this entire union in $\CP^{D-1}$.
Iteration will give the assertion.

For each ordered pair of components, consider pairs of distinct points
$(x,y)\in V_\alpha\times V_\beta$, deleting the diagonal when
$\alpha=\beta$. The two tautological lines representing $x$ and $y$
span a rank-two vector bundle over this pair space. Its projectivization
is a $\CP^1$-bundle whose evaluation map to $\CP^D$ sweeps out exactly
the secant lines joining those pairs. Its complex dimension is
\[
 \dim_{\C}V_\alpha+\dim_{\C}V_\beta+1\leq2d+1<D.
\]
The image has measure zero by Sard's theorem. The pair space need not
be compact; applying Sard in its countably many coordinate charts
still gives the same conclusion.

We also exclude every embedded projective tangent space. For a
component of complex dimension $e$, a local homogeneous lift and its
first derivatives span an $(e+1)$-dimensional vector space at each
point. These spaces are independent of the lift and coordinates and
form a vector bundle over that component. The evaluation map from its
projectivization sweeps out the embedded projective tangent spaces.
Its domain has complex dimension $2e\leq2d<D$, so this image too has
measure zero. The union $U$ itself has measure zero.

There are only finitely many components and component pairs. We may
therefore choose a point $q\in\CP^D$ outside $U$, all the secant lines
just described, and all the projective tangent spaces. Projection from
$q$ is holomorphic on $\CP^D\setminus\{q\}$. Two distinct points of
$U$ could have the same projection only if their secant line contained
$q$; hence the restricted projection is injective. The linear quotient
by the line representing $q$ is injective on each of the
$(e+1)$-dimensional tangent spans described above. Its projectivization
therefore has injective derivative on each component. The union is
compact, so this injective
holomorphic immersion is an embedding, with disjoint component images.
The same argument applies at each subsequent projection step, down to
$\CP^{2d+1}$.
\end{proof}

\subsection{The integral blow-up formula}

We now explain how the torsion of the centers passes to the ambient
manifold. An integral direct-sum formula is essential here: rational
cohomology alone would discard the information used in the final
comparison. The following classical formula holds without any
freeness assumption on the homology of the center; see
\cite[Theorem~4.1]{DuanLi2009}.

\begin{proposition}
\label{geo:blowup}
Let $A$ be a compact connected K\"ahler manifold and
$Z\subset A$ a nonempty closed connected complex submanifold of complex
codimension $c\geq2$. The complex blow-up
$\pi:\widehat A\to A$ is compact, connected, and K\"ahler.
If $A$ is simply connected, then $\widehat A$ is simply connected.
For every integer $i$,
\begin{equation}
\label{geo:blowup-homology}
 H_i(\widehat A;\Z)\cong
 H_i(A;\Z)\oplus\bigoplus_{j=1}^{c-1}H_{i-2j}(Z;\Z),
\end{equation}
where homology in negative degrees is zero.
\end{proposition}

\begin{proof}
Write $\nu$ for the complex normal bundle of $Z$ in $A$. Locally the
blow-up replaces a normal vector by a line together with a vector on
that line. It is a smooth complex manifold, and $\pi$ is a proper
holomorphic map that is an isomorphism away from $Z$. The exceptional
divisor is $E_Z=\mathbb P(\nu)$, with projection $\rho:E_Z\to Z$;
its normal line bundle is the tautological bundle $\mathcal O_{E_Z}(-1)$.
The local description gives connectedness, and properness over compact
$A$ gives compactness.

For K\"ahlerness, let $\omega_A$ be a K\"ahler form on $A$. A Hermitian
metric on $\nu$ induces a metric on $\mathcal O_{E_Z}(1)$ whose Chern
curvature is positive on each projective fiber. Extend this metric to
$\mathcal O(-E_Z)$, whose restriction to $E_Z$ is
$\mathcal O_{E_Z}(1)$, and let $\eta$ be its real Chern curvature form.
The form $\pi^*\omega_A$ is semipositive. It is positive away from
$E_Z$, and its kernel at a point of $E_Z$ consists precisely of the
tangent directions to the projective fiber. The form $\eta$ is positive
on those directions. On the unit tangent bundle for any auxiliary
metric, it remains positive in a neighborhood of the kernel directions;
on the compact complement, $\pi^*\omega_A$ has a positive lower bound
and $\eta$ is bounded. It follows that
\[
 \pi^*\omega_A+\varepsilon\eta
\]
is positive for all sufficiently small $\varepsilon>0$. It is closed
and of type $(1,1)$, so it is a K\"ahler form.

If $A$ is simply connected, general position in real codimension
$2c\geq4$ shows that $A\setminus Z$ is simply connected: paths and
contracting disks can be perturbed to avoid $Z$. Every loop in
$\widehat A$ can be perturbed off the real codimension two submanifold
$E_Z$, and then contracts in
$\widehat A\setminus E_Z\cong A\setminus Z$. Thus $\widehat A$ is
simply connected. No assumption on the fundamental group of the center
is needed.

We prove the splitting first in integral cohomology. Let
$\iota:E_Z\hookrightarrow\widehat A$ be the inclusion and put
$h=c_1(\mathcal O_{E_Z}(1))$. The integral Leray--Hirsch theorem gives
\begin{equation}
\label{geo:exceptional-cohomology}
 H^k(E_Z;\Z)=
 \bigoplus_{j=0}^{c-1}h^j\rho^*H^{k-2j}(Z;\Z).
\end{equation}
The classes $1,h,\ldots,h^{c-1}$ restrict to an integral basis on each
fiber; thus this statement does not require the base cohomology to be
free. The Gysin map $\iota_*$ raises degree by two and satisfies the
self-intersection identity
\[
 \iota^*\iota_*(v)=c_1(\mathcal O_{E_Z}(-1))v=-hv.
\]
Consequently the maps
\[
 s_j:H^{k-2j}(Z;\Z)\longrightarrow H^k(\widehat A;\Z),
 \qquad s_j(u)=-\iota_*(h^{j-1}\rho^*u),
 \qquad 1\leq j\leq c-1,
\]
lift all the positive powers of $h$ in
\eqref{geo:exceptional-cohomology}.

It remains to identify the classes whose restriction to $E_Z$ has
only an $h^0$ term. The map of pairs
$(\widehat A,E_Z)\to(A,Z)$ induces an isomorphism on relative
cohomology: collapsing the indicated subspaces gives a homeomorphism
$\widehat A/E_Z\cong A/Z$. Denote this relative cohomology isomorphism
by $\pi_{\mathrm{rel}}^*$, and let $\delta_A$ and $\delta_{\widehat A}$
be the connecting maps of the two pairs. If
$v\in H^k(\widehat A;\Z)$ satisfies $\iota^*v=\rho^*z$ for
$z\in H^k(Z;\Z)$, naturality and exactness give
\[
 \pi_{\mathrm{rel}}^*\delta_Az
 =\delta_{\widehat A}\rho^*z
 =\delta_{\widehat A}\iota^*v=0.
\]
The relative isomorphism implies $\delta_Az=0$, so $z$ extends to
$u\in H^k(A;\Z)$. The difference $v-\pi^*u$ restricts to zero on
$E_Z$, and hence comes from $H^k(\widehat A,E_Z;\Z)$. Its relative
preimage corresponds under $\pi_{\mathrm{rel}}^*$ to a class in
$H^k(A,Z;\Z)$. The absolute image of that class, added to $u$, pulls
back to $v$. Thus $v$ belongs to $\pi^*H^k(A;\Z)$.

Subtracting the lifts $s_j$ from an arbitrary cohomology class now shows
that
\[
 \pi^*+\sum_{j=1}^{c-1}s_j:
 H^k(A;\Z)\oplus\bigoplus_{j=1}^{c-1}H^{k-2j}(Z;\Z)
 \longrightarrow H^k(\widehat A;\Z)
\]
is surjective. It is injective as well. Restriction to $E_Z$ and
\eqref{geo:exceptional-cohomology} first force all coefficients of the
positive powers of $h$ to vanish. The remaining summand is detected
by $\pi^*$, which is injective over $\Z$: the map $\pi$ has degree one,
and the pushforward defined using integral Poincar\'e duality is a left
inverse to $\pi^*$. We have proved
\begin{equation}
\label{geo:blowup-cohomology}
 H^k(\widehat A;\Z)\cong
 H^k(A;\Z)\oplus\bigoplus_{j=1}^{c-1}H^{k-2j}(Z;\Z).
\end{equation}
If $\dim_{\C}A=n$, duality in degree $2n-i$ turns the center summands
into $H_{i-2c+2j}(Z;\Z)$. Replacing $j$ by $c-j$ gives
\eqref{geo:blowup-homology}. All maps and splittings were integral, so
the formula includes torsion as direct summands.
\end{proof}

\subsection{The family of ninefolds}

We apply the preceding results to the centers $Z_2$ and $Z_3$ from
\eqref{geo:centers}.

\begin{proposition}
\label{geo:family}
For every integer $m\geq1$ there is a simply connected closed K\"ahler
manifold $X_m$ of complex dimension nine, obtained by blowing up
$\CP^9$ along $m$ disjoint copies of each of $Z_2$ and $Z_3$. Its
integral homology is
\begin{equation}
\label{geo:family-homology}
\begin{split}
 H_i(X_m;\Z)\cong{}&H_i(\CP^9;\Z)\\
 &\oplus\left(\bigoplus_{r=1}^{6}H_{i-2r}(Z_2;\Z)\right)^{\oplus m}\\
 &\oplus\left(\bigoplus_{r=1}^{4}H_{i-2r}(Z_3;\Z)\right)^{\oplus m}.
\end{split}
\end{equation}
\end{proposition}

\begin{proof}
Apply Lemma~\ref{geo:projection} with $d=4$ to the disjoint union of
$m$ copies of $Z_2$ and $m$ copies of $Z_3$. It embeds all $2m$ centers
simultaneously in $\CP^9$. Blow up the centers successively. Because
they are disjoint, each remaining center and a neighborhood of it are
unchanged by the preceding blow-ups. Their complex codimensions are
always seven and five. Proposition~\ref{geo:blowup} preserves
compactness, connectedness, K\"ahlerness, and simple connectivity at
every step, starting from $\CP^9$; the complex dimension remains nine.
Applying \eqref{geo:blowup-homology} once for each center gives
\eqref{geo:family-homology}.
\end{proof}
